\documentclass[12pt]{article}
\usepackage[utf8]{inputenc}
\usepackage[T1]{fontenc}
\usepackage{hyperref}
\title{Uda Language Resources for K-12 Curriculum Development}
\author{Compiled by Claude Sonnet 5 (claude-sonnet-5) for Andr\'as Kornai}
\date{}
\begin{document}
\maketitle

\section{Introduction}
Uda (ISO 639-3 uda; Glottocode udaa1238) is a Lower Cross River language of Akwa Ibom State, Nigeria, spoken in Mbo Local Government Area. It is treated together with Enwang as Enwang-Uda. Estimates range from about 25,000 speakers (Wikipedia) to 26,000 people (Joshua Project). A Yale paper reports that nearly all speakers are also bilingual in Ibibio and that Uda persists mainly in home and village life. Very little material was found, and none of it was educational.

\section{Primary Educational Resources}
No primers, textbooks or school programmes for Uda were found in the searches.

\section{Linguistic and Reference Materials}
\begin{itemize}
\item Wikipedia entry on Enwang-Uda (classification, speaker numbers): \url{https://en.wikipedia.org/wiki/Enwang-Uda_language}
\item Overview of Lower Cross River languages: \url{https://en.wikipedia.org/wiki/Lower_Cross_River_languages}
\item Yale linguistics paper on endangered Uda and its bilingualism with Ibibio: \url{https://ling.yale.edu/media/330/download}
\end{itemize}
The 2012 CoLang field methods course at the University of Kansas worked on Uda. Any resulting materials were not located.

\section{Religious and Community Resources}
\begin{itemize}
\item Joshua Project profile of the Uda people (Mbo LGA): \url{https://joshuaproject.net/people_groups/15682}
\end{itemize}

\section{Digital and Audio Resources}
No Uda audio, video or software was found. For the related majority language, Ibibio, see Wikipedia (\url{https://en.wikipedia.org/wiki/Ibibio_language}) and Omniglot (\url{https://www.omniglot.com/writing/ibibio.htm}).

\section{Curriculum Development Context}
Documentation is thin and no orthography or teaching materials were found. Most speakers also speak Ibibio, the regional language, so a programme would probably need to be bilingual with Ibibio. The first steps are to find the CoLang 2012 data and to consult the community about an orthography.
\end{document}